\input{../preamble.tex}
\newgeometry{a4paper, margin=2cm, right=4cm}
\def\arraystretch{1}
\setlength\tabcolsep{5pt}

\SetDate[12/05/2026]

\newcommand{\exeframe}[4]{\exe{#1}{#2}{#3}{#4}\marginpar{\framebox(80,30)}}
\newcommand{\exemulticolsframe}[5]{\exemulticols{#1}{#2}{#3}{#4}{#5}\marginpar{\framebox(80,30)}}

\begin{document}
\pagestyle{fancy}
\fancyhead[L]{Seconde 9}
\fancyhead[C]{\textbf{Évaluation — Statistiques}}
\fancyhead[R]{\today}

\null\vspace{-20pt}
Consignes particulières : 
\begin{itemize}[label=$\bullet$]
	\item 
	La calculatrice est {autorisée}.
	\item
	Arrondir à $10^{-2}$ si nécessaire.
\end{itemize}

\hrule
\marginpar{\vspace{-5pt}\bf\underline{Réponses}}

\exeframe{}{
	Calculer la moyenne de la série statistique
		$ X = (1 ; 20 ; 17 ; 19 ; -20 ; -1). $ 
}{exe:1}{
	%
}

\exeframe{}{
	Donner la médiane de la série statistique
		\[ X = (1 ; 2 ; 4 ; 4 ; 6 ; 6 ; 7 ;7 ; 9 ; 9 ; 61). \]
}{exe:2}{
	%
}

\exemulticolsframe{}{
	Donner l'effectif total de la série statistique décrite par le diagramme à barres ci-contre.
}{
	\begin{center}
	\begin{tikzpicture}[scale=.8]
    \begin{axis}[
        ymin=0, ymax=3,
        xmin=0, xmax=20,
        minor x tick num = 1,
        %minor y tick num=1,
        axis line style = -,
        xtick distance=2,
        area style,
        xlabel = {Valeur},
        ylabel = {Effectif},
        xlabel near ticks,
        ylabel near ticks,
		extra x ticks={0},
		extra y ticks={0},
		x=10pt,
		ybar,
		bar width=8pt,
      ]
      \addplot+[bar shift=4pt,draw=BLACK,thick, fill=GREEN_E!50] plot coordinates {
        (2,1)
        (3,1)
        (5,2)
        (6,3)
        (7,1)
        (8,1)
        (9,1)
        (10,2)
        (11,1)
        (12,3)
        (13,2)
        (14,2)
        (15,2)
        (16,2)
        (17,1)
        (18,1)
        (19,1)
      };
    \end{axis} 
  \end{tikzpicture}
  \end{center}
\vspace{-40pt}\null
}{exe:3}{
	%
}

\exemulticolsframe{}{
	Calculer la moyenne de la série statistique ci-contre.
}{
	\begin{tabular}{|c|c|c|c|}\hline
	Valeur   & 0 & 20 & 10 \\ \hline
	Effectif & 17 & 17 & 1 \\ \hline
	\end{tabular}
}{exe:4}{
	%
}

\exemulticolsframe{}{
	Pour quel $N\in\N$ la moyenne de la série ci-contre vaut-elle 10 ?
}{
	\begin{tabular}{|c|c|c|c|}\hline
	Valeur   & 0 & 20 & 10 \\ \hline
	Effectif & 17 & 17 & 1 \\ \hline
	\end{tabular}
}{exe:5}{
	%
}

\exeframe{}{
	Un élève reçoit cinq notes listées ainsi : $(4,25 ; 10; 18,5 ; 15,5 ; 11,25)$.
	À ces notes sont associés cinq coefficients, donnés dans le même ordre par $(0,5 ; 1,4 ; 0,1 ; 0,5; 1,1)$.
	
	Calculer la moyenne pondérée des notes.
}{exe:6}{
	%
}

\exeframe{}{
	Quels coefficients donner aux valeurs de la série $X=(13 ; 18)$ pour que la moyenne pondérée soit égale à 15 ?
}{exe:8}{

}

\exeframe{}{
	Est-il possible de trouver deux coefficients tels que la moyenne pondérée de la série $X = (15 ; 18 ; 21 ; 35)$ soit égale à 40 ?
}{exe:9}{

}

\exeframe{}{
	Est-il possible de trouver deux coefficients tels que la moyenne pondérée de la série $X = (15 ; -1 ; 21 ; 35)$ soit égale à 0 ?
}{exe:10}{

}

\exeframe{}{
	La moyenne des notes d'une classe est donnée égale à $9,45$.
	
	Donner la nouvelle moyenne des notes après avoir ajouté $1,5$ points à chaque note.
}{exe:11}{

}

\exeframe{}{
	La moyenne des notes d'une classe est égale à $9,45$.
	
	Donner la nouvelle moyenne des notes après avoir ajouté $1,5$ points à chaque note.
}{exe:12}{

}

\exeframe{}{
	Est-ce que l'évolution réciproque d'une augmentation de 6\% est une diminution de 6\% ?
}{exe:13}{

}

\exeframe{}{
	Quel est le coefficient multiplicateur associé à une augmentation de 12\% ?
}{exe:14}{

}

\exeframe{}{
	La moyenne des salaires d'une entreprise est égale à 1930€.
	
	Donner la nouvelle moyenne des salaires après que chaque salaire ait été multiplié par $1,3$.
}{exe:15}{

}

\exeframe{}{
	À quel taux d'augmentation est associée la multiplication par $1,3$ ?
}{exe:16}{

}

\exeframe{}{
	Soient $\sigma_1$ et $\sigma_2$ les écarts-types de chaque série statistique suivante.
	A-t-on $\sigma_1 > \sigma_2$ ou $\sigma_1 < \sigma_2$ ?
	
	
	\begin{center}
	\begin{tikzpicture}[scale=.8]
    \begin{axis}[
        ymin=0, ymax=3,
        xmin=0, xmax=20,
        minor x tick num = 1,
        %minor y tick num=1,
        axis line style = -,
        xtick distance=2,
        area style,
        xlabel = {Valeur},
        ylabel = {Effectif},
        xlabel near ticks,
        ylabel near ticks,
		extra x ticks={0},
		extra y ticks={0},
		x=10pt,
		ybar,
		bar width=8pt,
      ]
      \addplot+[bar shift=4pt,draw=BLACK,thick, fill=GREEN_E!50] plot coordinates {
        (2,1)
        (3,1)
        (5,2)
        (6,3)
        (7,1)
        (8,1)
        (9,1)
        (10,2)
        (11,1)
        (12,3)
        (13,2)
        (14,2)
        (15,2)
        (16,2)
        (17,1)
        (18,1)
        (19,1)
      };
    \end{axis} 
  \end{tikzpicture}
	\begin{tikzpicture}[scale=.8]
    \begin{axis}[
        ymin=0, ymax=3,
        xmin=0, xmax=20,
        minor x tick num = 1,
        %minor y tick num=1,
        axis line style = -,
        xtick distance=2,
        area style,
        xlabel = {Valeur},
        ylabel = {Effectif},
        xlabel near ticks,
        ylabel near ticks,
		extra x ticks={0},
		extra y ticks={0},
		x=10pt,
		ybar,
		bar width=8pt,
      ]
      \addplot+[bar shift=4pt,draw=BLACK,thick, fill=GREEN_E!50] plot coordinates {
        (2,1)
        (3,1)
        (5,2)
        (6,3)
        (7,1)
        (8,1)
        (9,1)
        (10,2)
        (11,1)
        (12,3)
        (13,2)
        (14,2)
        (15,2)
        (16,2)
        (17,1)
        (18,1)
        (19,1)
      };
    \end{axis} 
  \end{tikzpicture}
  \end{center}
}{exe:7}{
	%
}


\exeframe{}{
	Calculer la moyenne de la série statistique suivante.

	\begin{center}
	\begin{tikzpicture}[scale=.8]
	\begin{axis}[
		xmin=0, xmax=14,
		ymin=0, ymax=8,
		x=20pt,
		y=10pt,
		grid=none,
		xtick distance=1,
		axis line style = -,
		extra x ticks={0},
		extra y ticks={0},
		xlabel=Valeur,
		xlabel near ticks,
		ylabel=Effectif,
		ylabel near ticks,
		ybar,
		bar width=8pt,
	]     
      \addplot+[bar shift=0,draw=BLACK, thick, fill=BLUE_E!50] plot coordinates {
        (3,8)
        (4,6)
        (5,2)
        (9,4)
        (10,3)
        (12,5)
        (13,2)
      };

	\end{axis}
	\end{tikzpicture}
	\end{center}
}{exe:17}{
	%
}


\exeframe{}{
	Donner le premier quartile de la série statistique suivante.

	\begin{center}
	\begin{tikzpicture}[scale=.8]
	\begin{axis}[
		xmin=0, xmax=14,
		ymin=0, ymax=8,
		x=20pt,
		y=10pt,
		grid=none,
		xtick distance=1,
		axis line style = -,
		extra x ticks={0},
		extra y ticks={0},
		xlabel=Valeur,
		xlabel near ticks,
		ylabel=Effectif,
		ylabel near ticks,
		ybar,
		bar width=8pt,
	]     
      \addplot+[bar shift=0,draw=BLACK, thick, fill=BLUE_E!50] plot coordinates {
        (3,8)
        (4,6)
        (5,2)
        (9,4)
        (10,3)
        (12,5)
        (13,2)
      };

	\end{axis}
	\end{tikzpicture}
	\end{center}
}{exe:18}{
	%
}

\exeframe{}{
	Donner le troisième quartile de la série statistique $X = (4 ; 5 ; 3 ; 1 ; 4 ; 2)$.
}{exe:19}{
	%
}

\exeframe{}{
	Donner l'étendue de la série statistique $X = (-3 ; 2,3 ; 4 ; 5 ; -1 ; 7)$.
}{exe:20}{
	%
}

%%%%%%%%%%%

\label{lastpage}
\newpage
\fancyhead[C]{\textbf{Solutions}}
\shipoutAnswer
	
\end{document}
